% Part: second-order-logic
% Chapter: sol-and-sets
% Section: cardinalities

\documentclass[../../../include/open-logic-section]{subfiles}

\begin{document}

\olfileid{sol}{set}{crd}

\olsection{సమితుల పరిమాణాలు}

\begin{explain}
వ్యక్తి క్షేత్రం పరిమితమా, అనంతమా,
\tetoken{లెక్కించదగినదా}{!!{enumerable}}, \tetoken{లెక్కించలేనిదా}{!!{nonenumerable}} అని
వ్యక్తం చేయగలిగినట్లే,~$\Domain{M}$ యొక్క ఒక ఉపసమితి పరిమితమా,
అనంతమా, \tetoken{లెక్కించదగినదా}{!!{enumerable}},
\tetoken{లెక్కించలేనిదా}{!!{nonenumerable}} అనే ధర్మాన్ని నిర్వచించవచ్చు.
\end{explain}

\begin{prop}
సూత్రం $\fn{Inf}(X) \ident$
\begin{multline*}
\lexists[u][(\lforall[x][(X(x) \lif X(u(x)))] \land {} \\
  \lforall[x][\lforall[y][((X(x) \land X(y) \land
      \eq[u(x)][u(y)]) \lif \eq[x][y])]] \land {} \\
  \lexists[y][(X(y) \land \lforall[x][(X(x)
      \lif \eq/[y][u(x)])])])]
\end{multline*}
చర నిర్దేశం~$s$కు సంబంధించి సంతృప్తిపరచబడితేనే $s(X)$ అనంతమైనది.
\sourcecorrection{OLTESOLSET-001}{ఆంగ్ల మూలంలోని $\fn{Inf}(X)$ సూత్రం $u$ను మొత్తం వ్యక్తి క్షేత్రంపై ఒకటి-ఒకటిగా మాత్రమే కోరింది; $u$ విలువలు $X$లోనే ఉండాలని గానీ, ఒకటి-ఒకటితనాన్ని $X$కు పరిమితం చేయాలని గానీ కోరలేదు. అందువల్ల అనంత వ్యక్తి క్షేత్రంలోని పరిమిత $X$ కూడా ఆ సూత్రాన్ని సంతృప్తిపరచగలదు. $u$ను $X$ నుంచి $X$కు ఒకటి-ఒకటి, సంగ్రస్తం కాని ప్రమేయంగా చెప్పే మూడు అవసరమైన నిబంధనలను ఇచ్చాం.}
\end{prop}

\begin{prop}
సూత్రం $\fn{Count}(X) \ident $
\begin{multline*}
(\lnot \lexists[x][X(x)] \lor
\lexists[z][\lexists[u][(X(z) \land
    \lforall[x][(X(x) \lif X(u(x)))] \land {} \\
    \lforall[Y][((Y \subseteq X \land Y(z) \land
      \lforall[x][(Y(x) \lif Y(u(x)))]) \lif X = Y)])]]))
\end{multline*}
చర నిర్దేశం~$s$కు సంబంధించి సంతృప్తిపరచబడితేనే $s(X)$
\tetoken{లెక్కించదగినది}{!!{enumerable}}.
\sourcecorrection{OLTESOLSET-002}{ఆంగ్ల మూలంలోని $\fn{Count}(X)$ సూత్రం శూన్య సమితిని విడిచిపెట్టింది; అంతేకాక ఆగమన నిబంధనలో $Y\subseteq X$ను కోరకపోవడంతో, మొత్తం వ్యక్తి క్షేత్రమే $Y$గా తీసుకుంటే అది $X=\Domain{M}$ను బలవంతం చేసింది. మూల నిర్వచనంలోని ముగింపు కుండలీకరణలు కూడా తారుమారయ్యాయి. శూన్య సందర్భాన్ని చేర్చి, మిగతా సందర్భంలో $Y$ను $X$ ఉపసమితులకు పరిమితం చేసి, సూత్రాన్ని సరిగా మూసాం.}
\end{prop}

\tetoken{లెక్కించలేని}{!!{nonenumerable}} సమితులు ఉన్నాయని, నిజానికి అనంత పరిమాణాల వేర్వేరు
స్థాయులు అనంతంగా ఉన్నాయని కాంటర్ సిద్ధాంతం నుంచి మనకు తెలుసు.
సమితి సిద్ధాంతం సమితుల పరిమాణాలకు సంబంధించిన మొత్తం అంకగణితాన్ని
అభివృద్ధి చేసి, సమితులకు అనంత కార్డినల్ సంఖ్యలను నిర్దేశిస్తుంది.
పరిమిత సమితుల పరిమాణాలను కొలిచే కార్డినల్ సంఖ్యలుగా సహజ సంఖ్యలు
పనిచేస్తాయి. \tetoken{అనంతంగా లెక్కించదగిన}{!!{denumerable}} సమితుల పరిమాణమే
మొదటి అనంత కార్డినాలిటీ; దాన్ని~$\aleph_0$ (“అలెఫ్-నాట్” లేదా
“అలెఫ్-సున్నా”) అంటారు. తరువాతి అనంత పరిమాణం~$\aleph_1$.
లెక్కించదగినది కాకుండా (అంటే, పరిమాణం~$\aleph_0$ కాకుండా) ఒక
సమితికి ఉండగల అతి చిన్న పరిమాణం అది. “$X$ పరిమాణం $\aleph_0$”ను
$\fn{Aleph}_0(X) \liff \fn{Inf}(X) \land \fn{Count}(X)$గా నిర్వచించవచ్చు.
$X$ యొక్క అన్ని నిజ ఉపసమితులు పరిమితమైనవి లేదా పరిమాణం~$\aleph_0$
గలవి అయి, అది స్వయంగా అనంతమైనదీ పరిమాణం~$\aleph_0$ గలది కానిదీ
అయితేనే దాని
పరిమాణం $\aleph_1$. కాబట్టి దీన్ని సూత్రం
$\fn{Aleph_1}(X) \ident \lforall[Y][((Y \subseteq X \land
  \lnot (X = Y)) \lif (\lnot\fn{Inf}(Y) \lor
  \fn{Aleph}_0(Y)))] \land \fn{Inf}(X) \land \lnot
\fn{Aleph}_0(X)$తో వ్యక్తం చేయవచ్చు.
\sourcecorrection{OLTESOLSET-007}{ఆంగ్ల మూలంలోని $\fn{Aleph_1}(X)$ సూత్రం $Y\subseteq X$లో $Y=X$ను కూడా అనుమతిస్తూనే $X$ పరిమాణం $\aleph_0$ కాదని మాత్రమే కోరింది. $Y=X$ను తీసుకుంటే మొదటి నిబంధన $X$ను పరిమితంగా బలవంతం చేస్తుంది; నిజానికి ప్రతి పరిమిత $X$ మూల సూత్రాన్ని సంతృప్తిపరుస్తుంది. పరిమాణం $\aleph_1$ను వ్యక్తం చేయడానికి $Y$ను నిజ ఉపసమితులకు పరిమితం చేసి, $X$ అనంతమైనదని చేర్చాం.}
పరిమాణం $\aleph_2$ కావడాన్ని ఇదే విధంగా నిర్వచిస్తాం; తదితరాలు.

ప్రత్యేక ఆసక్తి గల ఒక పరిమాణం ఉంది; దాన్ని అవిచ్ఛిన్న సమితి
కార్డినాలిటీ అంటారు. అది $\Pow{\Nat}$ పరిమాణం; సమానార్థకంగా,
అది~$\Real$ పరిమాణం. ఒక సమితి అవిచ్ఛిన్న సమితి పరిమాణం గలదని
ద్వితీయ-స్థాయి తర్కంలో కూడా వ్యక్తం చేయవచ్చు; కానీ అందుకు ఇంకొంత
పని అవసరం.

\end{document}
